\section{Proof of well-posedness}
\label{sec:wp-proof}

The route is: backward $L^2$ estimate
(Lemma~\ref{lem:bpbsde}) $\to$ small-coupling base case (Lemma~\ref{lem:small}) $\to$
uniform stability on $[\lambda_*,1]$ (Lemma~\ref{lem:pert}), which at $\lambda=1$ gives
Theorems~\ref{thm:stab} and~\ref{thm:uniq} $\to$ fixed-increment continuation
(Lemma~\ref{lem:cont}) $\to$ Theorem~\ref{thm:cwp}, the continuation being along the
coupling strength, in the tradition of Peng and Wu~\cite{PengWu}. The central a priori
estimate rests on a mixed $G$-energy identity for $\inner{G\,\delta X_t}{\delta Y_t}$,
whose compensated-Poisson cross-variation produces a coercive $\Lnu$ jump term controlled
by the $G$-monotonicity assumption. The continuation is in the coupling strength, not the
horizon: shortening the horizon does not make the frozen-coefficient map contractive,
since the $(z,u)$-dependence of $\sigma,h$ and the full-tuple law coupling enter the
contraction constant with no $T$-factor.

Inside $\diff t$-integrals we do not distinguish $X_t,Y_t$ from $X_{t-},Y_{t-}$:
c\`adl\`ag paths have only countably many jumps on compact intervals, so
$\Theta_t^-=\Theta_t$ $\Prob$-a.s.\ and $\mu_t^-=\mu_t$ for a.e.\ $t$
(Lemma~\ref{lem:lawflow}).

Throughout this section, $C$ denotes a positive constant, not necessarily the same at each
occurrence, depending only on $T,L,\|G\|,c_G,\alpha,\beta_1,\beta_2$ and the BDG constant
$c_{\mathrm{BDG}}$, unless a different dependence is stated; recurring constants are
subscripted ($C_0$, $C_{\mathrm{BP}}$, $C_B$, \dots) and fixed where they are introduced.

It\^o's formula for L\'evy-type stochastic integrals, and the product form obtained by
applying it to the bilinear map $(x,y)\mapsto\inner{Gx}{y}$, are those
of~\cite[\S4.4]{Applebaum}, stated for the Brownian--Poisson setting fixed above. The
maximal estimates use the Burkholder--Davis--Gundy inequality at $p=1$ in
quadratic-variation form (\cite{HeWangYan}, chapter on $H^1$ and BMO); valid for
c\`adl\`ag local martingales, it covers the compensated-Poisson martingales below without
a separate jump statement. By localisation and dominated convergence, a c\`adl\`ag local
martingale whose running supremum is integrable has zero-mean increments (the
\emph{supremum criterion} below).

For $\lambda\in[0,1]$ we work throughout with the \emph{$\lambda$-scaled, forced}
system: for a square-integrable forcing
$q=(q^b,q^\sigma,q^h,q^f,q^g_T)$ with
\[
\begin{aligned}
q^b&\in L^2(\Omega\times[0,T];\R^n),
&\qquad q^\sigma&\in L^2(\Omega\times[0,T];\R^{n\times d}),\\
q^h&\in L^2(\Omega\times[0,T]\times E;\R^n),
&\qquad q^f&\in L^2(\Omega\times[0,T];\R^m),\\
q^g_T&\in L^2(\Omega,\F_T;\R^m),
\end{aligned}
\]
where $q^b,q^\sigma,q^f$ are progressively
measurable and $q^h$ is predictable ($\mathcal P\otimes\mathcal E$-measurable as
in~\eqref{eq:fubiniiso}, square-integrable for $\diff\Prob\otimes\diff t\otimes\nu$),
\begin{equation}\label{eq:forced}
\left\{\;\begin{aligned}
X_t&=\xi_0+\lambda\!\int_0^t b(s,\Theta_s,\mu_s)\,\diff s
+\lambda\!\int_0^t\sigma(s,\Theta_s,\mu_s)\,\diff W_s\\
&\hspace{10.2em}+\lambda\!\int_0^t\!\!\int_E h(s,\Theta_s^-,e,\mu_s^-)\,\Nt(\diff s,\diff e)\\
&\hspace{5.4em}+\int_0^t q_s^b\,\diff s+\int_0^t q_s^\sigma\,\diff W_s
+\int_0^t\!\!\int_E q_s^h(e)\,\Nt(\diff s,\diff e),\\[2.2ex]
Y_t&=\lambda\,g(X_T,\Law(X_T))+\lambda\!\int_t^T f(s,\Theta_s,\mu_s)\,\diff s
+q_T^g+\int_t^T q_s^f\,\diff s\\
&\hspace{12.6em}-\int_t^T Z_s\,\diff W_s-\int_t^T\!\!\int_E U_s(e)\,\Nt(\diff s,\diff e).
\end{aligned}\right.
\end{equation}
Call the system \eqref{eq:forced} \emph{well posed at} $\lambda$ --- written
$\mathsf{WP}(\lambda)$ --- if it has a unique solution in $\HH$ for \emph{every} such
forcing $q$. The unforced target is
$\mathsf{WP}(1)$ with $q\equiv0$, so it suffices to prove that
$\mathsf{WP}(1)$ holds.

By Assumptions~\ref{ass:integ}--\ref{ass:lip}, the coefficient processes~\eqref{eq:coefmeas}
along any $\Theta\in\HH$ are square-integrable: for $\psi\in\{b,\sigma,h,f\}$, in the norms of
Assumption~\ref{ass:lip},
\begin{equation}\label{eq:coeffint}
\E\!\int_0^T\!|\psi(s,\Theta_s,\mu_s)|^2\,\diff s
\le C\,\E\!\int_0^T\!\Big(|X_s|^2+|Y_s|^2+|Z_s|_F^2+\Lnorm{U_s}^2
+|\psi(s,0,\delta_0)|^2\Big)\diff s,
\end{equation}
by that Lipschitz bound together with the coupling estimate \eqref{eq:coupling} at
$\mu'=\delta_0$. For the terminal function the analogous bound
\begin{equation}\label{eq:termg}
\E|g(X_T,\Law(X_T))|^2\le C\big(\E|X_T|^2+|g(0,\delta_0)|^2\big)
\end{equation}
follows in the same way from \eqref{eq:lipT} and \eqref{eq:couplingT}.

\subsection{A Brownian--Poisson \texorpdfstring{$L^2$}{L2} estimate}

\begin{lemma}[Brownian--Poisson $L^2$ estimate]\label{lem:bpbsde}
Let $\delta Y$ be a c\`adl\`ag adapted process and let $\delta Z\in\HW$, $\delta U\in\HN$
satisfy
\begin{equation}\label{eq:dYlin}
\delta Y_t=\delta\xi+\int_t^T\phi_s\,\diff s-\int_t^T\delta Z_s\,\diff W_s-\int_t^T\!\int_E\delta U_s(e)\,\Nt(\diff s,\diff e)
\end{equation}
with $\delta\xi\in L^2(\F_T)$ and $\phi\in L^2(\Omega\times[0,T])$. Then
$\delta Y\in\Ssp{\R^m}$, and there is a constant
$C_{\mathrm{BP}}$, depending only on $T$ and $c_{\mathrm{BDG}}$, with
\[
\E\sup_{t\le T}|\delta Y_t|^2+\E\!\int_0^T\!\big(|\delta Z_s|_F^2+\Lnorm{\delta U_s}^2\big)\diff s
\le C_{\mathrm{BP}}\Big(\E|\delta\xi|^2+\E\!\int_0^T|\phi_s|^2\diff s\Big).
\]
\end{lemma}
\begin{proof}
\emph{Membership.} By \eqref{eq:dYlin} with $\int_t^T=\int_0^T-\int_0^t$,
\[
\sup_{t\le T}|\delta Y_t|\le|\delta\xi|+\int_0^T\!|\phi_s|\,\diff s
+2\sup_{t\le T}\Big|\int_0^t\delta Z_s\,\diff W_s\Big|
+2\sup_{t\le T}\Big|\int_0^t\!\!\int_E\delta U_s(e)\,\Nt(\diff s,\diff e)\Big|.
\]
Both stochastic integrals are true $L^2$-martingales by the It\^o and
compensated-Poisson isometries, as $\delta Z\in\HW$ and $\delta U\in\HN$.
Squaring, Cauchy--Schwarz on the drift term, and Doob's inequality then give
$\E\sup_{t\le T}|\delta Y_t|^2<\infty$, that is, $\delta Y\in\Ssp{\R^m}$.

The $\diff s$- and $\diff W$-integrals of $\delta Y$ are continuous, so
$\delta Y$ jumps only through its compensated Poisson integral: its jumps
$\Delta\,\delta Y_s:=\delta Y_s-\delta Y_{s-}$ satisfy $\Delta\,\delta Y_s=\delta U_s(e)$
at each atom $(s,e)$ of $N$; hence the jump variation
splits through $N=\Nt+\diff s\,\nu$ as
\begin{align*}
\sum_{t<s\le T}|\Delta\,\delta Y_s|^2
&=\int_t^T\!\!\int_E|\delta U_s(e)|^2\,N(\diff s,\diff e)\\
&=\int_t^T\!\!\int_E|\delta U_s(e)|^2\,\Nt(\diff s,\diff e)
+\int_t^T\!\Lnorm{\delta U_s}^2\,\diff s.
\end{align*}
It\^o's formula for
$e^{t}|\delta Y_t|^2$ then gives the pathwise identity
\begin{equation}\label{eq:itoY2}
\begin{aligned}
e^{t}|\delta Y_t|^2
&+\int_t^T\!e^{s}\big(|\delta Y_s|^2+|\delta Z_s|_F^2+\Lnorm{\delta U_s}^2\big)\diff s\\
&=e^{T}|\delta\xi|^2+2\int_t^T\!e^{s}\inner{\delta Y_s}{\phi_s}\diff s\\
&\quad-2\int_t^T\!e^{s}\inner{\delta Y_{s-}}{\delta Z_s\,\diff W_s}\\
&\quad-2\int_t^T\!\!\int_E e^{s}\inner{\delta Y_{s-}}{\delta U_s(e)}\,\Nt(\diff s,\diff e)\\
&\quad-\int_t^T\!\!\int_E e^{s}|\delta U_s(e)|^2\,\Nt(\diff s,\diff e).
\end{aligned}
\end{equation}

\emph{Zero-mean terms.} The stochastic-integral terms in \eqref{eq:itoY2} decompose into three
local martingales
\[
\begin{aligned}
M^W_t&:=\int_0^t e^{s}\inner{\delta Y_{s-}}{\delta Z_s\,\diff W_s},\\
M^1_t&:=\int_0^t\!\int_E e^{s}\inner{\delta Y_{s-}}{\delta U_s(e)}\,\Nt(\diff s,\diff e),\\
M^2_t&:=\int_0^t\!\int_E e^{s}|\delta U_s(e)|^2\,\Nt(\diff s,\diff e),
\end{aligned}
\]
each of which we now show satisfies $\E[M_T-M_t]=0$. The three do not share one argument:
$\delta U\in\HN$ gives the quadratic integrand of $M^2$ no square moments, so the $L^2$
theory behind the bounds below does not reach it --- nor, with $\nu(E)=\infty$ admitted,
does localisation along jump times, which then fail to exhaust $[0,T]$. $M^2$ is instead
handled through the compensation formula, which asks only for integrability.
For $M^W$, BDG and Cauchy--Schwarz give
\[
\E\sup_{t\le T}|M^W_t|
\le c_{\mathrm{BDG}}\,e^{T}\big(\E\sup_{t\le T}|\delta Y_t|^2\big)^{1/2}
\Big(\E\!\int_0^T\!|\delta Z_s|_F^2\,\diff s\Big)^{1/2}<\infty ,
\]
and for $M^1$ BDG, Cauchy--Schwarz and \eqref{eq:compid} give
\begin{align*}
\E\sup_{t\le T}|M^1_t|
&\le c_{\mathrm{BDG}}\,\E\Big[\Big(\int_0^T\!\!\int_E e^{2s}|\delta Y_{s-}|^2|\delta U_s(e)|^2
\,N(\diff s,\diff e)\Big)^{1/2}\Big]\\
&\le c_{\mathrm{BDG}}\,e^{T}\big(\E\sup_{t\le T}|\delta Y_t|^2\big)^{1/2}
\Big(\E\!\int_0^T\!\!\int_E|\delta U_s(e)|^2\,N(\diff s,\diff e)\Big)^{1/2}\\
&=c_{\mathrm{BDG}}\,e^{T}\big(\E\sup_{t\le T}|\delta Y_t|^2\big)^{1/2}
\Big(\E\!\int_0^T\!\!\int_E|\delta U_s(e)|^2\,\nu(\diff e)\,\diff s\Big)^{1/2}\\
&=c_{\mathrm{BDG}}\,e^{T}\big(\E\sup_{t\le T}|\delta Y_t|^2\big)^{1/2}
\Big(\E\!\int_0^T\!\Lnorm{\delta U_s}^2\,\diff s\Big)^{1/2}<\infty .
\end{align*}
These bounds put the running suprema of $M^W$ and $M^1$ in $L^1$, so the supremum
criterion gives $\E[M_T-M_t]=0$ for $M\in\{M^W,M^1\}$.
For $M^2$ the integrand $e^{s}|\delta U_s(e)|^2$ is predictable and nonnegative and,
because $\delta U\in\HN$, lies in $L^1(\diff\Prob\otimes\diff s\otimes\nu)$:
\[
\E\!\int_0^T\!\!\int_E e^{s}|\delta U_s(e)|^2\,\nu(\diff e)\,\diff s
\le e^{T}\,\E\!\int_0^T\!\Lnorm{\delta U_s}^2\,\diff s<\infty .
\]
Both the Poisson integral and its compensator are therefore integrable, so the increment
splits under the expectation:
\begin{align*}
\E[M^2_T-M^2_t]
&=\E\!\int_t^T\!\!\int_E e^{s}|\delta U_s(e)|^2\,\Nt(\diff s,\diff e)\\
&=\E\!\int_t^T\!\!\int_E e^{s}|\delta U_s(e)|^2\,N(\diff s,\diff e)
-\E\!\int_t^T\!\!\int_E e^{s}|\delta U_s(e)|^2\,\nu(\diff e)\,\diff s\\
&=0 .
\end{align*}
Taking
expectations in \eqref{eq:itoY2} therefore removes the three martingale terms, and
bounding the drift term by Young's inequality,
$2\inner{\delta Y_s}{\phi_s}\le|\delta Y_s|^2+|\phi_s|^2$, cancels the $|\delta Y_s|^2$
integrals on the two sides:
\[
\E e^{t}|\delta Y_t|^2
+\E\int_t^T e^{s}\big(|\delta Z_s|_F^2+\Lnorm{\delta U_s}^2\big)\diff s
\le
\E e^{T}|\delta\xi|^2
+\E\int_t^T e^{s}|\phi_s|^2\diff s .
\]
Bounding $1\le e^{s}\le e^{T}$ on both sides leaves, for every $t\le T$,
\[
\E|\delta Y_t|^2+\E\int_t^T\big(|\delta Z_s|_F^2+\Lnorm{\delta U_s}^2\big)\diff s\le R ,
\qquad
R:=e^{T}\Big(\E|\delta\xi|^2+\E\int_0^T|\phi_s|^2\diff s\Big).
\]
Both terms on the left are nonnegative and $R$ does not depend on $t$, so each is bounded
by $R$ on its own, the first uniformly in $t$ and the second at $t=0$; adding them,
\begin{equation}\label{eq:supE}
\begin{aligned}
&\sup_{t\le T}\E|\delta Y_t|^2
+\E\int_0^T\big(|\delta Z_s|_F^2+\Lnorm{\delta U_s}^2\big)\diff s\\
&\qquad\le 2R
=C_0\Big(\E|\delta\xi|^2+\E\int_0^T|\phi_s|^2\diff s\Big),
\qquad
C_0:=2e^{T} .
\end{aligned}
\end{equation}

It remains to upgrade $\sup_{t\le T}\E|\delta Y_t|^2$ to $\E\sup_{t\le T}|\delta Y_t|^2$. Taking the difference of
\eqref{eq:dYlin} at times $t$ and $0$ gives the forward form
\[
\delta Y_t=\delta Y_0-\int_0^t\phi_s\diff s+\int_0^t\delta Z_s\diff W_s+\int_0^t\!\int_E\delta U_s(e)\Nt(\diff s,\diff e) .
\]
Squaring, taking the supremum over $t$, Cauchy--Schwarz on the $\diff s$-integral, and BDG on
the two martingales give
\begin{align*}
\E\sup_{t\le T}|\delta Y_t|^2
&\le 4\,\E|\delta Y_0|^2+4\,\E\Big(\int_0^T\!|\phi_s|\diff s\Big)^{2}\\
&\qquad+4\,\E\sup_{t\le T}\Big|\int_0^t\!\delta Z_s\diff W_s\Big|^{2}
+4\,\E\sup_{t\le T}\Big|\int_0^t\!\!\int_E\delta U_s(e)\,\Nt(\diff s,\diff e)\Big|^{2}\\
&\le 4\,\E|\delta Y_0|^2+4T\,\E\!\int_0^T\!|\phi_s|^2\diff s
+4c_{\mathrm{BDG}}\,\E\!\int_0^T\!\big(|\delta Z_s|_F^2+\Lnorm{\delta U_s}^2\big)\diff s .
\end{align*}
Using $\E|\delta Y_0|^2\le\sup_{t\le T}\E|\delta Y_t|^2$ and \eqref{eq:supE},
\begin{align*}
\E\sup_{t\le T}|\delta Y_t|^2
&+\E\!\int_0^T\!\big(|\delta Z_s|_F^2+\Lnorm{\delta U_s}^2\big)\diff s\\
&\le 4\sup_{t\le T}\E|\delta Y_t|^2
+(4c_{\mathrm{BDG}}+1)\,\E\!\int_0^T\!\big(|\delta Z_s|_F^2+\Lnorm{\delta U_s}^2\big)\diff s\\
&\qquad+4T\,\E\!\int_0^T\!|\phi_s|^2\diff s\\
&\le\big(4C_0+(4c_{\mathrm{BDG}}+1)C_0+4T\big)
\Big(\E|\delta\xi|^2+\E\!\int_0^T\!|\phi_s|^2\diff s\Big)\\
&=C_{\mathrm{BP}}\Big(\E|\delta\xi|^2+\E\!\int_0^T\!|\phi_s|^2\diff s\Big),
\qquad C_{\mathrm{BP}}:=C_0(5+4c_{\mathrm{BDG}})+4T .
\end{align*}
This depends only on $T$ and $c_{\mathrm{BDG}}$.
\end{proof}

\subsection{Small-coupling solvability}

\begin{lemma}[Small-coupling base case]\label{lem:small}
Under Assumptions~\ref{ass:integ}--\ref{ass:lip}, there is $\lambda_*\in(0,1]$, depending only on $T$, $L$ and
$c_{\mathrm{BDG}}$, such that $\mathsf{WP}(\lambda)$ holds for
every $\lambda\in[0,\lambda_*]$.
\end{lemma}

\begin{proof}
Fix $\lambda\in[0,1]$ and a forcing $q$. For a trial tuple $\Theta=(X,Y,Z,U)\in\HH$ write
\[
\mu^\Theta_t:=\Law(\Theta_t),\qquad\mu^{\Theta,-}_t:=\Law(\Theta_t^-),
\]
and define $\Phi_\lambda(\Theta):=\bar\Theta=(\bar X,\bar Y,\bar Z,\bar U)$ by the forward and
backward constructions below.

\emph{Forward.} Let $\bar X$ be given by the frozen forward dynamics
\begin{multline}\label{eq:frozenfwd}
\bar X_t=\xi_0+\lambda\!\int_0^t b(s,\Theta_s,\mu^\Theta_s)\,\diff s
+\lambda\!\int_0^t\sigma(s,\Theta_s,\mu^\Theta_s)\,\diff W_s\\
+\lambda\!\int_0^t\!\!\int_E h(s,\Theta_s^-,e,\mu^{\Theta,-}_s)\,\Nt(\diff s,\diff e)\\
+\int_0^t q^b_s\,\diff s+\int_0^t q^\sigma_s\,\diff W_s+\int_0^t\!\!\int_E q^h_s(e)\,\Nt(\diff s,\diff e).
\end{multline}
In \eqref{eq:frozenfwd} the coefficient $h$ is evaluated at the left limits
$(\Theta^-_s,\mu^{\Theta,-}_s)$; by Lemma~\ref{lem:lawflow} these agree with
$(\Theta_s,\mu^\Theta_s)$ outside a $\diff t\otimes\diff\Prob$-null set, so either form may be
used inside the $\diff t\otimes\diff\Prob$-integrals below.
Since $\Theta\in\HH$, \eqref{eq:coeffint} gives
\[
\E\!\int_0^T\!|\psi(s,\Theta_s,\mu^\Theta_s)|^2\,\diff s<\infty
\qquad\text{for }\psi\in\{b,\sigma,h\},
\]
and $q^b,q^\sigma,q^h$ are square-integrable by hypothesis. The coefficients are jointly Borel,
so these integrands are progressively measurable and the jump integrand is predictable.
Cauchy--Schwarz on the drift and Doob/BDG on the two martingale integrals then give
$\E\sup_{t\le T}|\bar X_t|^2<\infty$, so \eqref{eq:frozenfwd} defines a unique
$\bar X\in\Ssp{\R^n}$.

\emph{Backward.} The random variable
\[
\Xi:=\lambda\,g(\bar X_T,\Law(\bar X_T))+q^g_T+\int_0^T(\lambda f(s,\Theta_s,\mu^\Theta_s)+q^f_s)\,\diff s
\]
lies in $L^2(\Omega,\F_T;\R^m)$ by \eqref{eq:termg}, \eqref{eq:coeffint} applied to $f$, and the
square-integrability of $q^g_T$ and $q^f$. Set $M_t:=\E[\Xi\mid\F_t]$. By the Brownian--Poisson
predictable representation property there are unique
$\bar Z\in\HW$, $\bar U\in\HN$ with
\[
M_t=M_0+\int_0^t\bar Z_s\,\diff W_s+\int_0^t\!\int_E\bar U_s(e)\,\Nt(\diff s,\diff e).
\]
Put
\[
\bar Y_t:=M_t-\int_0^t(\lambda f(s,\Theta_s,\mu^\Theta_s)+q^f_s)\,\diff s .
\]
Doob's inequality and $\Xi\in L^2$ give $M\in\Ssp{\R^m}$. The drift is square-integrable by
\eqref{eq:coeffint} and $q^f\in L^2$, so Cauchy--Schwarz gives $\bar Y\in\Ssp{\R^m}$. By construction
$(\bar Y,\bar Z,\bar U)$ solves the frozen linear BSDE
\begin{multline}\label{eq:frozenbwd}
\bar Y_t=\lambda g(\bar X_T,\Law(\bar X_T))+q^g_T
+\int_t^T\!(\lambda f(s,\Theta_s,\mu^\Theta_s)+q^f_s)\,\diff s\\
-\int_t^T\bar Z_s\,\diff W_s-\int_t^T\!\!\int_E\bar U_s(e)\,\Nt(\diff s,\diff e),
\end{multline}
so $\bar\Theta\in\HH$ and $\Phi_\lambda:\HH\to\HH$ is well defined.

\emph{Contraction.} Take $\Theta^1,\Theta^2\in\HH$, write $\bar\Theta^k=\Phi_\lambda(\Theta^k)$,
and set
\[
\begin{gathered}
\delta\psi_s:=\psi(s,\Theta^1_s,\mu^{\Theta^1}_s)-\psi(s,\Theta^2_s,\mu^{\Theta^2}_s)
\quad(\psi\in\{b,\sigma,h,f\}),\\[0.4ex]
\delta X:=X^1-X^2,\qquad \delta\bar X:=\bar X^1-\bar X^2 ,
\end{gathered}
\]
and likewise for the remaining components of $\Theta^1-\Theta^2$ and of
$\bar\Theta^1-\bar\Theta^2$; the common initial datum $\xi_0$ and
the fixed forcing $q$ cancel in every difference.

\emph{Step 1: the coefficient differences.} Write
$S_s:=|\delta b_s|+|\delta\sigma_s|_F+\Lnorm{\delta h_s}+|\delta f_s|$. Assumption~\ref{ass:lip},
applied to the diagonal inputs $(\Theta^1_s,\mu^{\Theta^1}_s)$ and $(\Theta^2_s,\mu^{\Theta^2}_s)$,
gives
\[
S_s\le L\Big(|\delta X_s|+|\delta Y_s|+|\delta Z_s|_F+\Lnorm{\delta U_s}
+W_2\big(\mu^{\Theta^1}_s,\mu^{\Theta^2}_s\big)\Big) .
\]
Squaring, using Cauchy--Schwarz in $\R^5$, taking expectations, and bounding the deterministic
Wasserstein term by \eqref{eq:coupling} give
\[
\E\,S_s^2\le 10L^2\,\E\big[|\delta X_s|^2+|\delta Y_s|^2+|\delta Z_s|_F^2+\Lnorm{\delta U_s}^2\big].
\]
Integrating in $s$ and dominating the $X$- and $Y$-slots by their suprema,
\begin{equation}\label{eq:coeffdiff}
\begin{aligned}
\E\!\int_0^T\! S_s^2\,\diff s
&\le 10L^2\,\E\!\int_0^T\!\big(|\delta X_s|^2+|\delta Y_s|^2+|\delta Z_s|_F^2+\Lnorm{\delta U_s}^2\big)\diff s\\
&\le 10L^2\Big(T\,\E\sup_{t\le T}|\delta X_t|^2+T\,\E\sup_{t\le T}|\delta Y_t|^2\\
&\hspace{6.2em}+\E\!\int_0^T\!\big(|\delta Z_s|_F^2+\Lnorm{\delta U_s}^2\big)\diff s\Big)\\
&\le L_\star^2\,\|\Theta^1-\Theta^2\|_{\HH}^2 ,
\qquad L_\star^2:=10L^2\max\{T,1\} .
\end{aligned}
\end{equation}

\emph{Step 2: the forward difference.} Taking the difference of \eqref{eq:frozenfwd} at
$\Theta^1$ and $\Theta^2$,
\[
\delta\bar X_t=\lambda\!\int_0^t\!\delta b_s\,\diff s+\lambda\!\int_0^t\!\delta\sigma_s\,\diff W_s
+\lambda\!\int_0^t\!\!\int_E\delta h_s(e)\,\Nt(\diff s,\diff e) .
\]
Squaring, taking the supremum in $t$, and taking expectations,
\begin{equation}\label{eq:fwddiff}
\begin{aligned}
\E\sup_{t\le T}|\delta\bar X_t|^2
&\le 3\lambda^2\Big[\E\Big(\int_0^T\!|\delta b_s|\,\diff s\Big)^{2}
+\E\sup_{t\le T}\Big|\int_0^t\!\delta\sigma_s\,\diff W_s\Big|^{2}\\
&\hspace{13.4em}+\E\sup_{t\le T}\Big|\int_0^t\!\!\int_E\delta h_s(e)\,\Nt(\diff s,\diff e)\Big|^{2}\Big]\\
&\le 3\lambda^2\Big[T\,\E\!\int_0^T\!|\delta b_s|^2\,\diff s
+c_{\mathrm{BDG}}\,\E\!\int_0^T\!|\delta\sigma_s|_F^2\,\diff s
+c_{\mathrm{BDG}}\,\E\!\int_0^T\!\Lnorm{\delta h_s}^2\,\diff s\Big]\\
&\le 3\lambda^2\max\{T,c_{\mathrm{BDG}}\}\,
\E\!\int_0^T\!\big(|\delta b_s|^2+|\delta\sigma_s|_F^2+\Lnorm{\delta h_s}^2\big)\diff s\\
&\le 3\lambda^2\max\{T,c_{\mathrm{BDG}}\}\,\E\!\int_0^T\! S_s^2\,\diff s\\
&\le C_X\lambda^2\,\|\Theta^1-\Theta^2\|_{\HH}^2 ,
\qquad C_X:=3\max\{T,c_{\mathrm{BDG}}\}\,L_\star^2 ,
\end{aligned}
\end{equation}
the jump term by \eqref{eq:compid}, the next step since the summands of $S_s$ are
non-negative, and the last by \eqref{eq:coeffdiff}.

\emph{Step 3: the terminal difference.} Put
\[
\delta\bar\xi:=\lambda\big[g(\bar X^1_T,\Law(\bar X^1_T))-g(\bar X^2_T,\Law(\bar X^2_T))\big].
\]
Then
\begin{equation}\label{eq:terdiff}
\begin{aligned}
\E|\delta\bar\xi|^2
&\le\lambda^2L^2\,\E\big(|\delta\bar X_T|+W_2(\Law(\bar X^1_T),\Law(\bar X^2_T))\big)^2\\
&\le2\lambda^2L^2\big(\E|\delta\bar X_T|^2+W_2(\Law(\bar X^1_T),\Law(\bar X^2_T))^2\big)\\
&\le C_g\lambda^2\,\E|\delta\bar X_T|^2 ,\qquad C_g:=4L^2 ,
\end{aligned}
\end{equation}
the first step by \eqref{eq:lipT} and the last by \eqref{eq:couplingT}.

\emph{Step 4: the backward difference.} Taking the difference of \eqref{eq:frozenbwd} at
$\Theta^1$ and $\Theta^2$, the terms $q^g_T$ and $q^f$ cancelling,
\[
\delta\bar Y_t=\delta\bar\xi+\int_t^T\!\lambda\,\delta f_s\,\diff s
-\int_t^T\!\delta\bar Z_s\,\diff W_s-\int_t^T\!\!\int_E\delta\bar U_s(e)\,\Nt(\diff s,\diff e) ,
\]
which is \eqref{eq:dYlin} with $\delta\xi=\delta\bar\xi$ and $\phi_s=\lambda\,\delta f_s$.
$\delta\bar Y$ is c\`adl\`ag adapted, and the remaining hypotheses hold:
\[
\begin{alignedat}{2}
(\delta\bar Z,\delta\bar U)&\in\HW\times\HN
&\quad&\text{since }\bar\Theta^1,\bar\Theta^2\in\HH,\\
\delta\bar\xi&\in L^2(\F_T) &&\text{by }\eqref{eq:termg},\\
\lambda\,\delta f&\in L^2(\Omega\times[0,T]) &&\text{by }\eqref{eq:coeffdiff},
\text{ as }|\delta f|\le S.
\end{alignedat}
\]
By Lemma~\ref{lem:bpbsde},
\begin{equation}\label{eq:bwddiff}
\begin{aligned}
\E\sup_{t\le T}|\delta\bar Y_t|^2
&+\E\!\int_0^T\!\big(|\delta\bar Z_s|_F^2+\Lnorm{\delta\bar U_s}^2\big)\diff s\\
&\le C_{\mathrm{BP}}\Big(\E|\delta\bar\xi|^2+\lambda^2\,\E\!\int_0^T\!|\delta f_s|^2\,\diff s\Big)\\
&\le C_{\mathrm{BP}}\Big(C_g\lambda^2\,\E|\delta\bar X_T|^2
+L_\star^2\lambda^2\,\|\Theta^1-\Theta^2\|_{\HH}^2\Big)\\
&\le C_{\mathrm{BP}}\Big(C_g\lambda^2\,\E\sup_{t\le T}|\delta\bar X_t|^2
+L_\star^2\lambda^2\,\|\Theta^1-\Theta^2\|_{\HH}^2\Big)\\
&\le C_{\mathrm{BP}}\big(C_gC_X\lambda^4+L_\star^2\lambda^2\big)\|\Theta^1-\Theta^2\|_{\HH}^2\\
&\le C_Y\lambda^2\,\|\Theta^1-\Theta^2\|_{\HH}^2 ,
\qquad C_Y:=C_{\mathrm{BP}}\big(C_gC_X+L_\star^2\big) ,
\end{aligned}
\end{equation}
the second inequality by \eqref{eq:terdiff} and \eqref{eq:coeffdiff}, the fourth by
\eqref{eq:fwddiff}, and the last by $\lambda^4\le\lambda^2$ for $\lambda\in[0,1]$.

\emph{Step 5: the contraction.} Since
$\Phi_\lambda(\Theta^1)-\Phi_\lambda(\Theta^2)
=(\delta\bar X,\delta\bar Y,\delta\bar Z,\delta\bar U)$,
\begin{equation}\label{eq:contract}
\begin{aligned}
\|\Phi_\lambda(\Theta^1)-\Phi_\lambda(\Theta^2)\|_{\HH}^2
&=\E\sup_{t\le T}|\delta\bar X_t|^2\\
&\quad+\E\sup_{t\le T}|\delta\bar Y_t|^2
+\E\!\int_0^T\!\big(|\delta\bar Z_s|_F^2+\Lnorm{\delta\bar U_s}^2\big)\diff s\\
&\le(C_X+C_Y)\lambda^2\,\|\Theta^1-\Theta^2\|_{\HH}^2\\
&\le C_\Phi\lambda^2\,\|\Theta^1-\Theta^2\|_{\HH}^2 ,
\qquad C_\Phi:=\max\{1,C_X+C_Y\} ,
\end{aligned}
\end{equation}
the first inequality by \eqref{eq:fwddiff} and \eqref{eq:bwddiff}. Here
\[
L_\star=L_\star(T,L),\quad C_X=C_X(T,L,c_{\mathrm{BDG}}),\quad
C_g=C_g(L),\quad C_{\mathrm{BP}}=C_{\mathrm{BP}}(T,c_{\mathrm{BDG}}),
\]
so $C_\Phi=C_\Phi(T,L,c_{\mathrm{BDG}})$. Setting
\[
\lambda_*:=(2\sqrt{C_\Phi})^{-1}\in(0,\tfrac12]
\]
gives $C_\Phi\lambda^2\le\tfrac14$ for $\lambda\in[0,\lambda_*]$, so that \eqref{eq:contract}
reads
\[
\|\Phi_\lambda(\Theta^1)-\Phi_\lambda(\Theta^2)\|_{\HH}
\le\tfrac12\,\|\Theta^1-\Theta^2\|_{\HH},\qquad\lambda\in[0,\lambda_*] .
\]
By the Banach fixed-point theorem, $\Phi_\lambda$ has a unique fixed point $\Theta\in\HH$.
Evaluating $\Phi_\lambda$ at $\Theta$ itself,
\[
\mu^\Theta=\mu,\qquad\bar\Theta=\Theta ,
\]
so \eqref{eq:frozenfwd}--\eqref{eq:frozenbwd} become the forward and backward equations of
\eqref{eq:forced}, which is solved by $\Theta$.

\emph{Conversely.} Let $\Theta\in\HH$ solve \eqref{eq:forced} and evaluate $\Phi_\lambda$
at the trial $\Theta$. Again $\mu^\Theta=\mu$, so \eqref{eq:frozenfwd} and the forward
equation of \eqref{eq:forced} have the same right-hand side, which does not involve
$\bar X$; hence $\bar X=X$. For the backward component set
\[
M'_t:=Y_t+\int_0^t\!\big(\lambda f(s,\Theta_s,\mu_s)+q^f_s\big)\,\diff s ,
\]
so that
\[
M'_T=\Xi,\qquad
\diff M'_t=Z_t\,\diff W_t+\int_E U_t(e)\,\Nt(\diff t,\diff e) .
\]
Since $Z\in\HW$ and $U\in\HN$, $M'$ is a true $L^2$-martingale, so
\[
M'_t=\E[\Xi\mid\F_t]=M_t .
\]
Uniqueness of the predictable representation forces $(\bar Z,\bar U)=(Z,U)$, and
\eqref{eq:frozenbwd} then gives $\bar Y=Y$, so
\[
\bar\Theta=\Theta .
\]
Solutions are therefore exactly the fixed points, and the solution is unique; as $q$ was
arbitrary, $\mathsf{WP}(\lambda)$ holds on $[0,\lambda_*]$.
\end{proof}

\subsection{Continuous perturbation stability}

\begin{lemma}[Continuous perturbation stability along the homotopy]\label{lem:pert}
Let Assumptions~\ref{ass:lip}--\ref{ass:mono} hold, fix $\underline\lambda\in(0,1]$, and
fix $\lambda\in[\underline\lambda,1]$. Let
$\Theta^1,\Theta^2$ solve \eqref{eq:forced} at parameter $\lambda$ with $\F_0$-measurable
initial data $\xi_0^1,\xi_0^2$ and forcings
$q^k=(q^{b,k},q^{\sigma,k},q^{h,k},q^{f,k},q^{g,k}_T)$, $k=1,2$, where, for each $k$,
$X^k,Y^k$ are c\`adl\`ag adapted with
$\E\int_0^T(|X^k_t|^2+|Y^k_t|^2)\,\diff t<\infty$ and $X^k_T\in L^2(\Omega,\F_T;\R^n)$,
$Z^k\in\HW$, $U^k\in\HN$, and $\Theta^1-\Theta^2\in\HH$. Then, writing
$\delta\,{\cdot}={\cdot}^1-{\cdot}^2$,
\begin{multline}\label{eq:contpert}
\E\sup_{t\le T}|\delta X_t|^2+\E\sup_{t\le T}|\delta Y_t|^2
+\E\!\int_0^T\!\big(|\delta Z_t|_F^2+\Lnorm{\delta U_t}^2\big)\diff t\\
\le C_B\Big(\E|\delta\xi_0|^2+\E|\delta q^g_T|^2
+\E\!\int_0^T\!\big(|\delta q^b_t|^2+|\delta q^\sigma_t|_F^2
+\Lnorm{\delta q^h_t}^2+|\delta q^f_t|^2\big)\diff t\Big),
\end{multline}
where $C_B$ depends only on
$T,L,\|G\|,c_G,\alpha,\beta_1,\beta_2,\underline\lambda$ and $c_{\mathrm{BDG}}$.
\end{lemma}
\begin{proof}
The difference equations below involve the forcings only through $\delta q$, which we
therefore write simply as $q$. We define
\[
P:=\E|\delta\xi_0|^2+\E|q^g_T|^2
+\E\!\int_0^T\!\big(|q^b_t|^2+|q^\sigma_t|_F^2+\Lnorm{q^h_t}^2+|q^f_t|^2\big)\diff t .
\]
Set
\[
\delta\psi_t:=\psi(t,\Theta^1_t,\mu^1_t)-\psi(t,\Theta^2_t,\mu^2_t),\qquad
\psi\in\{b,\sigma,h,f\},
\]
where $\mu^k_t:=\Law(\Theta^k_t)$; likewise $\delta g_T$, with $\delta h$ evaluated at the
predictable tuples $(\Theta^{k,-}_t,\mu^{k,-}_t)$. For a.e.\ $t$ each pair
$(\Theta^k_t,\mu^k_t)$ is a diagonal input in the sense of Definition~\ref{def:diag}:
$\Theta^k_t\in L^2(\Omega,\F_t;\mathcal H)$ for a.e.\ $t$ by Fubini, and $\mu^k_t$ is
the Borel law flow of Lemma~\ref{lem:lawflow}.
The difference equations are
\begin{align}
\diff\,\delta X_t&=(\lambda\delta b_t+q^b_t)\,\diff t+(\lambda\delta\sigma_t+q^\sigma_t)\,\diff W_t
+\int_E(\lambda\delta h_t(e)+q^h_t(e))\,\Nt(\diff t,\diff e),\label{eq:dX}\\
\diff\,\delta Y_t&=-(\lambda\delta f_t+q^f_t)\,\diff t+\delta Z_t\,\diff W_t
+\int_E\delta U_t(e)\,\Nt(\diff t,\diff e),\qquad \delta Y_T=\lambda\,\delta g_T+q^g_T.\label{eq:dY}
\end{align}

The five steps run as follows. Steps~1--2 pair $G\,\delta X$ against $\delta Y$ and
insert Assumption~\ref{ass:mono}: this controls the projections $G\delta X$ and
$\Gstar\delta Y,\Gstar\delta Z,\Gstar\delta U$, but only by the data $P$ together with a
free multiple of the full norm. Steps~3--4 record two a~priori estimates, the forward one
bounding $\delta X$ through $(\delta Y,\delta Z,\delta U)$ and the backward one bounding
$(\delta Y,\delta Z,\delta U)$ through $\delta X$; neither closes alone.
Step~5 breaks the circle: full rank makes $\Gstar$ injective
when $n\ge m$ and $G$ injective when $m\ge n$, so the coercive blocks present shed their
projections, and which blocks are present decides whether the forward or the backward
estimate enters first. The free multiple is then chosen small and absorbed.

\emph{Step 1 ($G$-energy identity).} Both $\delta X$ and $\delta Y$ jump only through
their compensated Poisson integrals, with $\Delta\,\delta X_t=(\lambda\delta h_t+q^h_t)(e)$
and $\Delta\,\delta Y_t=\delta U_t(e)$ at each atom $(t,e)$ of $N$; hence the jump
covariation splits through $N=\Nt+\diff t\,\nu$ as
\[
\begin{aligned}
\sum_{t\le T}\inner{G\,\Delta\delta X_t}{\Delta\delta Y_t}
&=\int_0^T\!\!\int_E\inner{G(\lambda\delta h_t+q^h_t)(e)}{\delta U_t(e)}\,N(\diff t,\diff e)\\
&=\int_0^T\!\!\int_E\inner{G(\lambda\delta h_t+q^h_t)(e)}{\delta U_t(e)}\,\Nt(\diff t,\diff e)\\
&+\int_0^T\!\!\int_E\inner{\lambda\delta h_t(e)+q^h_t(e)}{\Gstar\delta U_t(e)}\,\nu(\diff e)\diff t .
\end{aligned}
\]
The continuous covariation of $\delta X$ and $\delta Y$, whose diffusions are
$\lambda\delta\sigma_t+q^\sigma_t$ and $\delta Z_t$, contributes
\[
\int_0^T\inner{\lambda\delta\sigma_t+q^\sigma_t}{\Gstar\delta Z_t}_F\,\diff t .
\]
It\^o's product formula for $\inner{G\,\delta X_t}{\delta Y_t}$, integrated over $[0,T]$ and using
$\delta X_0=\delta\xi_0$, then gives the pathwise identity
\begin{equation}\label{eq:Gid-path}
\begin{aligned}
\inner{G\delta X_T}{\delta Y_T}-\inner{G\,\delta\xi_0}{\delta Y_0}
&=\int_0^T\inner{\lambda\delta b_t+q^b_t}{\Gstar\delta Y_{t-}}\,\diff t\\
&+\int_0^T\inner{\lambda\delta\sigma_t+q^\sigma_t}{\Gstar\delta Z_t}_F\,\diff t\\
&+\int_0^T\!\!\int_E\inner{\lambda\delta h_t(e)+q^h_t(e)}{\Gstar\delta U_t(e)}\,\nu(\diff e)\diff t\\
&+\int_0^T\inner{-\lambda\delta f_t-q^f_t}{G\delta X_{t-}}\,\diff t\\
&+\int_0^T\inner{G\delta X_{t-}}{\delta Z_t\,\diff W_t}\\
&+\int_0^T\inner{G(\lambda\delta\sigma_t+q^\sigma_t)\,\diff W_t}{\delta Y_{t-}}\\
&+\int_0^T\!\!\int_E\inner{G\delta X_{t-}}{\delta U_t(e)}\,\Nt(\diff t,\diff e)\\
&+\int_0^T\!\!\int_E\inner{G(\lambda\delta h_t+q^h_t)(e)}{\delta Y_{t-}}\,\Nt(\diff t,\diff e)\\
&+\int_0^T\!\!\int_E\inner{G(\lambda\delta h_t+q^h_t)(e)}{\delta U_t(e)}\,\Nt(\diff t,\diff e) .
\end{aligned}
\end{equation}
Of the last five terms of \eqref{eq:Gid-path}, the first four are bounded by BDG and
Cauchy--Schwarz, with \eqref{eq:compid} in the jump terms:
\begin{equation}\label{eq:bdg4}
\left\{\;
\begin{aligned}
&\E\sup_{t\le T}\Big|\int_0^t\inner{G\delta X_{s-}}{\delta Z_s\,\diff W_s}\Big|\\
&\hspace{0.14\linewidth}\le c_{\mathrm{BDG}}\|G\|\big(\E\sup_{t\le T}|\delta X_t|^2\big)^{1/2}
\Big(\E\!\int_0^T\!|\delta Z_t|_F^2\,\diff t\Big)^{1/2},\\
&\E\sup_{t\le T}\Big|\int_0^t\inner{G(\lambda\delta\sigma_s+q^\sigma_s)\,\diff W_s}{\delta Y_{s-}}\Big|\\
&\hspace{0.14\linewidth}\le c_{\mathrm{BDG}}\|G\|\big(\E\sup_{t\le T}|\delta Y_t|^2\big)^{1/2}
\Big(\E\!\int_0^T\!|\lambda\delta\sigma_t+q^\sigma_t|_F^2\,\diff t\Big)^{1/2},\\
&\E\sup_{t\le T}\Big|\int_0^t\!\!\int_E\inner{G\delta X_{s-}}{\delta U_s(e)}\,\Nt(\diff s,\diff e)\Big|\\
&\hspace{0.14\linewidth}\le c_{\mathrm{BDG}}\|G\|\big(\E\sup_{t\le T}|\delta X_t|^2\big)^{1/2}
\Big(\E\!\int_0^T\!\Lnorm{\delta U_t}^2\,\diff t\Big)^{1/2},\\
&\E\sup_{t\le T}\Big|\int_0^t\!\!\int_E\inner{G(\lambda\delta h_s+q^h_s)(e)}{\delta Y_{s-}}\,\Nt(\diff s,\diff e)\Big|\\
&\hspace{0.14\linewidth}\le c_{\mathrm{BDG}}\|G\|\big(\E\sup_{t\le T}|\delta Y_t|^2\big)^{1/2}
\Big(\E\!\int_0^T\!\Lnorm{\lambda\delta h_t+q^h_t}^2\,\diff t\Big)^{1/2} .
\end{aligned}
\right.
\end{equation}

Since $\Theta^1-\Theta^2\in\HH$,
\[
\E\sup_{t\le T}|\delta X_t|^2+\E\sup_{t\le T}|\delta Y_t|^2
+\E\!\int_0^T\!\big(|\delta Z_t|_F^2+\Lnorm{\delta U_t}^2\big)\diff t
=\|\Theta^1-\Theta^2\|_{\HH}^2<\infty ,
\]
and by the computation of \eqref{eq:coeffdiff}, with the forcings square-integrable,
\[
\E\!\int_0^T\!\big(|\lambda\delta\sigma_t+q^\sigma_t|_F^2
+\Lnorm{\lambda\delta h_t+q^h_t}^2\big)\diff t<\infty .
\]
The four suprema in \eqref{eq:bdg4} are therefore integrable, so the supremum criterion
applies to those four local martingales. The fifth has predictable
integrand, and by Cauchy--Schwarz
\[
\begin{aligned}
&\E\!\int_0^T\!\!\int_E\big|\inner{\lambda\delta h_t(e)+q^h_t(e)}{\Gstar\delta U_t(e)}\big|\,\nu(\diff e)\diff t\\
&\hspace{0.26\linewidth}\le\|G\|\Big(\E\!\int_0^T\!\Lnorm{\lambda\delta h_t+q^h_t}^2\diff t\Big)^{1/2}
\Big(\E\!\int_0^T\!\Lnorm{\delta U_t}^2\diff t\Big)^{1/2}\\
&\hspace{0.26\linewidth}<\infty .
\end{aligned}
\]
Hence \eqref{eq:zeromean} applies, and all five terms have zero expectation.

Taking expectations and dropping the left limits in the drift multipliers (the
$\diff t$-integral convention above) reduces \eqref{eq:Gid-path} to
\begin{multline}\label{eq:Gid}
\E\inner{G\delta X_T}{\delta Y_T}-\E\inner{G\,\delta\xi_0}{\delta Y_0}
=\E\!\int_0^T\!\Big[\inner{\lambda\delta b_t+q^b_t}{\Gstar\delta Y_t}
+\inner{\lambda\delta\sigma_t+q^\sigma_t}{\Gstar\delta Z_t}_F\\
+\!\int_E\!\inner{\lambda\delta h_t(e)+q^h_t(e)}{\Gstar\delta U_t(e)}\nu(\diff e)
+\inner{-\lambda\delta f_t-q^f_t}{G\delta X_t}\Big]\diff t .
\end{multline}

\emph{Step 2 (insertion of monotonicity).} In the $h$-channel $\delta h_t$ is evaluated at the
predictable tuples $(\Theta^{k,-}_t,\mu^{k,-}_t)$, which agree with $(\Theta^k_t,\mu^k_t)$
for $\diff t$-a.e.\ $t$ by Lemma~\ref{lem:lawflow}. Assumption~\ref{ass:mono}, applied at
the common input pair $(\Theta^1_t,\mu^1_t),(\Theta^2_t,\mu^2_t)$, therefore bounds the
unperturbed part of the integrand of \eqref{eq:Gid}:
\begin{equation}\label{eq:intmono}
\begin{aligned}
&\lambda\E\Big[\inner{\delta b_t}{\Gstar\delta Y_t}+\inner{\delta\sigma_t}{\Gstar\delta Z_t}_F
+\!\int_E\!\inner{\delta h_t(e)}{\Gstar\delta U_t(e)}\nu(\diff e)+\inner{-\delta f_t}{G\delta X_t}\Big]\\
&\qquad\le-\lambda\beta_1\E\big[|\Gstar\delta Y_t|^2+|\Gstar\delta Z_t|_F^2+\Lnorm{\Gstar\delta U_t}^2\big]
-\lambda\beta_2\E|G\delta X_t|^2 .
\end{aligned}
\end{equation}

The terminal side of \eqref{eq:Gid} is bounded by the terminal monotonicity of
Assumption~\ref{ass:mono}:
\begin{equation}\label{eq:termmono}
\begin{aligned}
\E\inner{G\delta X_T}{\delta Y_T}
&=\lambda\E\inner{G\delta X_T}{\delta g_T}+\E\inner{G\delta X_T}{q^g_T}\\
&\ge\lambda\alpha\E|G\delta X_T|^2-|\E\inner{G\delta X_T}{q^g_T}| .
\end{aligned}
\end{equation}
Recall
\[
P=\E|\delta\xi_0|^2+\E|q^g_T|^2
+\E\!\int_0^T\!\big(|q^b_t|^2+|q^\sigma_t|_F^2+\Lnorm{q^h_t}^2+|q^f_t|^2\big)\diff t ,
\]
and set
\[
\begin{aligned}
\mathcal N'&:=\E|\delta X_T|^2+\E|\delta Y_0|^2
+\E\!\int_0^T\!\big(|\delta X_t|^2+|\delta Y_t|^2+|\delta Z_t|_F^2+\Lnorm{\delta U_t}^2\big)\diff t ,\\
\mathcal K&:=\eta\,\mathcal N'+\tfrac{\|G\|^2}{4\eta}\,P ,\qquad \eta>0 .
\end{aligned}
\]
Inserting \eqref{eq:intmono} and \eqref{eq:termmono} into \eqref{eq:Gid} gives
\[
\begin{aligned}
&\lambda\alpha\E|G\delta X_T|^2-|\E\inner{G\delta X_T}{q^g_T}|-\E\inner{G\,\delta\xi_0}{\delta Y_0}\\
&\qquad\le-\lambda\beta_1\E\!\int_0^T\!\big[|\Gstar\delta Y_t|^2+|\Gstar\delta Z_t|_F^2+\Lnorm{\Gstar\delta U_t}^2\big]\diff t
-\lambda\beta_2\E\!\int_0^T\!|G\delta X_t|^2\diff t\\
&\qquad+\E\!\int_0^T\!\Big(\inner{q^b_t}{\Gstar\delta Y_t}+\inner{q^\sigma_t}{\Gstar\delta Z_t}_F
+\inner{-q^f_t}{G\delta X_t}\\
&\qquad\qquad\quad+\int_E\inner{q^h_t(e)}{\Gstar\delta U_t(e)}\,\nu(\diff e)\Big)\diff t .
\end{aligned}
\]
Rearranging, and applying Young's inequality with $|\Gstar v|\le\|G\|\,|v|$ and
$|Gv|\le\|G\|\,|v|$, gives
\begin{equation}\label{eq:coerc}
\begin{aligned}
&\lambda\alpha\E|G\delta X_T|^2+\lambda\beta_2\E\!\int_0^T\!|G\delta X_t|^2\diff t\\
&\quad+\lambda\beta_1\E\!\int_0^T\!\big[|\Gstar\delta Y_t|^2+|\Gstar\delta Z_t|_F^2+\Lnorm{\Gstar\delta U_t}^2\big]\diff t\\
&\qquad\le|\E\inner{G\,\delta\xi_0}{\delta Y_0}|+|\E\inner{G\delta X_T}{q^g_T}|
+\E\!\int_0^T\!\big(|\inner{q^b}{\Gstar\delta Y}|+|\inner{q^\sigma}{\Gstar\delta Z}_F|\big)\diff t\\
&\qquad+\E\!\int_0^T\!\Bigl(\Bigl|\int_E\inner{q^h(e)}{\Gstar\delta U(e)}\,\nu(\diff e)\Bigr|
+|\inner{q^f}{G\delta X}|\Bigr)\diff t\\
&\qquad\le\mathcal K .
\end{aligned}
\end{equation}
Each coercive block present on the left of \eqref{eq:coerc} is therefore at most
$\mathcal K$.

\emph{Step 3 (forward estimate).} Fix $t\le T$ and integrate \eqref{eq:dX} over $[0,s]$,
$s\le t$, with $\delta X_0=\delta\xi_0$. Squaring, taking the supremum over $s$ and
expectations, Cauchy--Schwarz on the drift, BDG on the two martingales and
Assumption~\ref{ass:lip} give
\[
\begin{aligned}
\E\sup_{s\le t}|\delta X_s|^2
&\le C\Big(\E|\delta\xi_0|^2+\E\!\int_0^t\!\big(|\delta b_s|^2+|\delta\sigma_s|_F^2+\Lnorm{\delta h_s}^2\big)\diff s+P\Big)\\
&\le C\Big(\E|\delta\xi_0|^2+\E\!\int_0^t\!\big(|\delta X_s|^2+|\delta Y_s|^2+|\delta Z_s|_F^2+\Lnorm{\delta U_s}^2\big)\diff s+P\Big)\\
&\le C\Big(\int_0^t\!\E\sup_{r\le s}|\delta X_r|^2\,\diff s
+\E\!\int_0^T\!\big(|\delta Y_s|^2+|\delta Z_s|_F^2+\Lnorm{\delta U_s}^2\big)\diff s+P\Big) .
\end{aligned}
\]
The map $t\mapsto\E\sup_{s\le t}|\delta X_s|^2$ is nondecreasing and finite, since
$\Theta^1-\Theta^2\in\HH$, so Gr\"onwall gives
\begin{equation}\label{eq:fwd}
\E\sup_{t\le T}|\delta X_t|^2\le C\Big(\E\!\int_0^T\!\big(|\delta Y_t|^2+|\delta Z_t|_F^2+\Lnorm{\delta U_t}^2\big)\diff t+P\Big).
\end{equation}

\emph{Step 4 (backward estimate).} It\^o's formula applied to $e^{\kappa t}|\delta Y_t|^2$
gives, for $\kappa>0$,
\begin{equation}\label{eq:bwdito}
\begin{aligned}
&\E\,e^{\kappa t}|\delta Y_t|^2
+\E\!\int_t^T\!e^{\kappa s}\big(\kappa|\delta Y_s|^2+|\delta Z_s|_F^2+\Lnorm{\delta U_s}^2\big)\diff s\\
&\qquad\le\E\,e^{\kappa T}|\lambda\delta g_T+q^g_T|^2
+2\,\E\!\int_t^T\!e^{\kappa s}\inner{\delta Y_s}{\lambda\delta f_s+q^f_s}\diff s .
\end{aligned}
\end{equation}
By Assumption~\ref{ass:lip}, the coupling bound \eqref{eq:coupling}, and Young's
inequality, for every $\varepsilon>0$,
\begin{equation}\label{eq:bwdyoung}
2\,\E\inner{\delta Y_s}{\lambda\delta f_s+q^f_s}
\le \frac C\varepsilon\,\E|\delta Y_s|^2+\varepsilon\,\E\big(|\delta Z_s|_F^2+\Lnorm{\delta U_s}^2\big)
+\frac C\varepsilon\,\E\big(|\delta X_s|^2+|q^f_s|^2\big),
\end{equation}
with $C$ independent of $\varepsilon$.

Insert \eqref{eq:bwdyoung} in \eqref{eq:bwdito}:
\[
\begin{aligned}
&\E\,e^{\kappa t}|\delta Y_t|^2
+\E\!\int_t^T\!e^{\kappa s}\Big[\big(\kappa-\tfrac{C}{\varepsilon}\big)|\delta Y_s|^2
+(1-\varepsilon)\big(|\delta Z_s|_F^2+\Lnorm{\delta U_s}^2\big)\Big]\diff s\\
&\qquad\le\E\,e^{\kappa T}|\lambda\delta g_T+q^g_T|^2
+\tfrac{C}{\varepsilon}\,\E\!\int_t^T\!e^{\kappa s}\big(|\delta X_s|^2+|q^f_s|^2\big)\diff s .
\end{aligned}
\]
Take $\varepsilon<1$ and $\kappa>C/\varepsilon$. The $\delta Y$-integral on the left is then
nonnegative and may be dropped, and the weights satisfy $1\le e^{\kappa s}\le e^{\kappa T}$
on $[0,T]$, so
\[
\E|\delta Y_t|^2+\E\!\int_t^T\!\big(|\delta Z_s|_F^2+\Lnorm{\delta U_s}^2\big)\diff s
\le C\Big(\E|\lambda\delta g_T+q^g_T|^2
+\E\!\int_t^T\!\big(|\delta X_s|^2+|q^f_s|^2\big)\diff s\Big).
\]
With $\E|\delta g_T|^2\le C\,\E|\delta X_T|^2$, the forcings in $P$, and
$\int_t^T\le\int_0^T$ on the right, the bound no longer depends on $t$. Taking the
supremum over $t$ in the first term on the left and $t=0$ in the second, there remains
\begin{equation}\label{eq:wgt}
\sup_{t\le T}\E|\delta Y_t|^2
+\E\!\int_0^T\!\big(|\delta Z_t|_F^2+\Lnorm{\delta U_t}^2\big)\diff t
\le C\Big(\E|\delta X_T|^2+\E\!\int_0^T\!|\delta X_t|^2\diff t+P\Big).
\end{equation}
Writing \eqref{eq:dY} forward from $\delta Y_0$, BDG, Cauchy--Schwarz,
Assumption~\ref{ass:lip} with the coupling bound \eqref{eq:coupling}, and \eqref{eq:wgt} give
\[
\begin{aligned}
\E\sup_{t\le T}|\delta Y_t|^2
&\le C\Big(\E|\delta Y_0|^2+\E\!\int_0^T\!|\lambda\,\delta f_s+q^f_s|^2\diff s
+\E\!\int_0^T\!\big(|\delta Z_s|_F^2+\Lnorm{\delta U_s}^2\big)\diff s\Big)\\
&\le C\Big(\E|\delta Y_0|^2
+\E\!\int_0^T\!\big(|\delta X_s|^2+|\delta Y_s|^2+|\delta Z_s|_F^2+\Lnorm{\delta U_s}^2\big)\diff s+P\Big)\\
&\le C\Big(\E|\delta X_T|^2+\E\!\int_0^T\!|\delta X_t|^2\diff t+P\Big) .
\end{aligned}
\]
The first term of \eqref{eq:wgt} may therefore be replaced by
$\E\sup_{t\le T}|\delta Y_t|^2$:
\begin{equation}\label{eq:bwd}
\E\sup_{t\le T}|\delta Y_t|^2+\E\!\int_0^T\!\big(|\delta Z_t|_F^2+\Lnorm{\delta U_t}^2\big)\diff t
\le C\Big(\E|\delta X_T|^2+\E\!\int_0^T\!|\delta X_t|^2\diff t+P\Big).
\end{equation}

\emph{Step 5 (rank-case closure, then absorb).} Write
\[
\mathcal N:=\E\sup_{t\le T}|\delta X_t|^2+\E\sup_{t\le T}|\delta Y_t|^2+\E\!\int_0^T\!\big(|\delta Z_t|_F^2+\Lnorm{\delta U_t}^2\big)\diff t
\]
so that \eqref{eq:contpert} is the claim $\mathcal N\le C_B P$. Then
$\mathcal N'\le(1+2T)\mathcal N$, and by
\eqref{eq:coerc} each coercive block on its left is at most $\mathcal K$.

\emph{Case A ($n\ge m$; active when $m<n$, where $\beta_1>0$).} As $\Gstar$ is injective,
Assumption~\ref{ass:G} gives
\[
|\delta Y|^2\le c_G|\Gstar\delta Y|^2,\qquad
|\delta Z|_F^2\le c_G|\Gstar\delta Z|_F^2,\qquad
\Lnorm{\delta U}^2\le c_G\Lnorm{\Gstar\delta U}^2 .
\]
With $\beta_1>0$ and $\lambda\ge\underline\lambda$, the $\beta_1$-block of \eqref{eq:coerc} yields
\[
\E\!\int_0^T\!\big(|\delta Y_t|^2+|\delta Z_t|_F^2+\Lnorm{\delta U_t}^2\big)\diff t\le\bigl(c_G/(\underline\lambda\beta_1)\bigr)\mathcal K.
\]
Inserting this in \eqref{eq:fwd}, and the result in \eqref{eq:bwd},
\[
\begin{aligned}
\E\sup_{t\le T}|\delta X_t|^2&\le C(\mathcal K+P),\\
\E\sup_{t\le T}|\delta Y_t|^2
+\E\!\int_0^T\!\big(|\delta Z_t|_F^2+\Lnorm{\delta U_t}^2\big)\diff t&\le C(\mathcal K+P),
\end{aligned}
\]
so $\mathcal N\le C(\mathcal K+P)$.

\emph{Case B ($m\ge n$; active when $m>n$, where $\alpha,\beta_2>0$).} As $G$ is injective,
Assumption~\ref{ass:G} gives
\[
|\delta X|^2\le c_G|G\delta X|^2 .
\]
With $\alpha,\beta_2>0$ and $\lambda\ge\underline\lambda$, the $\alpha,\beta_2$-blocks of
\eqref{eq:coerc} yield
\[
\E|\delta X_T|^2+\E\!\int_0^T\!|\delta X_t|^2\diff t\le\bigl(2c_G/(\underline\lambda\min\{\alpha,\beta_2\})\bigr)\mathcal K.
\]
Inserting this in \eqref{eq:bwd}, and the result in \eqref{eq:fwd},
\[
\begin{aligned}
\E\sup_{t\le T}|\delta Y_t|^2
+\E\!\int_0^T\!\big(|\delta Z_t|_F^2+\Lnorm{\delta U_t}^2\big)\diff t&\le C(\mathcal K+P),\\
\E\sup_{t\le T}|\delta X_t|^2&\le C(\mathcal K+P),
\end{aligned}
\]
so again $\mathcal N\le C(\mathcal K+P)$.

When $m=n$, both $G,\Gstar$ are injective, so Case~A is available whenever $\beta_1>0$ and
Case~B whenever $\alpha,\beta_2>0$; Assumption~\ref{ass:mono} provides at least one of the
two, and when both are available we use whichever gives the smaller bound. In every
case
\[
\mathcal N\le C\Big(\eta\mathcal N'+\tfrac{\|G\|^2}{4\eta} P+P\Big)\le C\eta(1+2T)\mathcal N+C\,P;
\]
choosing $\eta$
small enough that $C\eta(1+2T)\le\tfrac12$ absorbs the $\mathcal N$-term (finite,
since $\Theta^1-\Theta^2\in\HH$) and gives
$\mathcal N\le C\,P$, which is \eqref{eq:contpert} with $C_B:=C$.
\end{proof}

\subsection{Stability}\label{sub:stabuniq}

\begin{proof}[Proof of Theorem~\ref{thm:stab}]
Take the system-$1$ coefficients $(b^1,\sigma^1,h^1,f^1,g^1)$ as the baseline and read
$\Theta^2$ as a forced solution of that system: $\Theta^2$ solves \eqref{eq:forced} at
$\lambda=1$ with
\[
q^{\psi,2}:=-(\delta\psi),\qquad \psi\in\{b,\sigma,h,f\},\qquad
q^{g,2}_T:=-\widehat{\delta g}_T ,
\]
while $\Theta^1$ solves it with zero forcing.

These form a square-integrable forcing for \eqref{eq:forced}: as differences of the
coefficient processes \eqref{eq:coefmeas}, the $q^{\psi,2}$ inherit their measurability,
and they are square-integrable by \eqref{eq:coeffint} applied to each system with its own
constants; $q^{g,2}_T\in L^2(\Omega,\F_T;\R^m)$ by \eqref{eq:lipT} with $X^2_T\in L^2$.

Since $q^1=0$, the difference $\delta q=q^1-q^2$ has $\delta q^\psi=(\delta\psi)$ and
$\delta q^g_T=\widehat{\delta g}_T$, so Lemma~\ref{lem:pert} at parameter $\lambda=1$ and
lower bound $\underline\lambda=1$, applied to
system~$1$, gives
\begin{multline*}
\E\sup_{t\le T}|\delta X_t|^2+\E\sup_{t\le T}|\delta Y_t|^2
+\E\!\int_0^T\!\big(|\delta Z_t|_F^2+\Lnorm{\delta U_t}^2\big)\diff t\\
\le C_B\Big(\E|\delta\xi_0|^2+\E\big|\widehat{\delta g}_T\big|^2\\
+\E\!\int_0^T\!\big(|(\delta b)(t)|^2+|(\delta\sigma)(t)|_F^2
+\Lnorm{(\delta h)(t)}^2+|(\delta f)(t)|^2\big)\diff t\Big),
\end{multline*}
which is \eqref{eq:stab} with $C_{\mathrm{stab}}:=C_B$.
\end{proof}

\subsection{The continuation step}

Taking identical sources and the common initial datum in Lemma~\ref{lem:pert} (at
lower bound $\underline\lambda=\lambda_*$) yields
uniqueness at each $\lambda\in[\lambda_*,1]$; Lemma~\ref{lem:small} gives it on
$[0,\lambda_*]$. It remains to prove existence up to $\lambda=1$.

\begin{lemma}[Continuation]\label{lem:cont}
Let Assumptions~\ref{ass:integ}--\ref{ass:mono} hold, and let $\lambda_*$ be the
small-coupling threshold of Lemma~\ref{lem:small}. There is $\zeta_0>0$, depending only on
$T$, $L$, $\|G\|$, $c_G$, $\alpha$, $\beta_1$, $\beta_2$, $\lambda_*$ and $c_{\mathrm{BDG}}$, such that: if
$\lambda_0\in[\lambda_*,1]$ and $\mathsf{WP}(\lambda_0)$ holds, then
$\mathsf{WP}(\lambda_0+\zeta)$ holds for every $\zeta\in[0,\zeta_0]$ with $\lambda_0+\zeta\le1$.
\end{lemma}

\begin{proof}
Fix $\lambda_0$ with $\mathsf{WP}(\lambda_0)$, set $\lambda_1=\lambda_0+\zeta\le1$, and
fix a forcing $q$ for the $\lambda_1$-system. For a trial $\Theta\in\HH$ define the
augmented forcing
\[
R^\psi(\Theta):=q^\psi+\zeta\,\psi(\cdot,\Theta,\mu^\Theta),\ \ \psi\in\{b,\sigma,h,f\},
\qquad R^g_T(\Theta):=q^g_T+\zeta\,g(X_T,\Law(X_T)).
\]
Then $R(\Theta)$ is
again a square-integrable forcing for \eqref{eq:forced}. Its integrability follows from
\eqref{eq:coeffint} and the Lipschitz/coupling bounds, and its $h$-channel is evaluated at
the predictable tuple $(\Theta^-,\mu^{\Theta,-})$, hence predictable.
Since $\mathsf{WP}(\lambda_0)$ holds, let $\Phi(\Theta):=\bar\Theta$ be the unique
solution in $\HH$ of the forced $\lambda_0$-system with forcing $R(\Theta)$.

For $\Theta^1,\Theta^2$, the outputs $\Phi(\Theta^k)$ solve the same $\lambda_0$-system
with sources $R(\Theta^1),R(\Theta^2)$; $q$ cancels in the difference, leaving
\[
R^\psi(\Theta^1)-R^\psi(\Theta^2)=\zeta[\psi(\cdot,\Theta^1,\mu^{\Theta^1})-\psi(\cdot,\Theta^2,\mu^{\Theta^2})].
\]
By Lemma~\ref{lem:pert} at parameter $\lambda_0$ and lower bound
$\underline\lambda=\lambda_*$, and the
Lipschitz bound of Assumption~\ref{ass:lip},
\[
\begin{aligned}
\|\Phi(\Theta^1)-\Phi(\Theta^2)\|_{\HH}^2
&\le C_B\!\sum_\psi\E\!\int_0^T|R^\psi(\Theta^1)-R^\psi(\Theta^2)|^2\diff t
+C_B\,\E|R^g_T(\Theta^1)-R^g_T(\Theta^2)|^2\\
&\le C_B\,C\,\zeta^2\,\|\Theta^1-\Theta^2\|_{\HH}^2 ,\qquad C=C(T,L) .
\end{aligned}
\]
Put $\zeta_0:=\min\{1,(2\sqrt{C_BC})^{-1}\}$, so that for $\zeta\le\zeta_0$
\[
\|\Phi(\Theta^1)-\Phi(\Theta^2)\|_{\HH}\le\tfrac12\,\|\Theta^1-\Theta^2\|_{\HH} .
\]
By the Banach fixed-point theorem $\Phi$ has a unique fixed point $\Theta$, at which
$R^\psi(\Theta)=q^\psi+\zeta\,\psi(\cdot,\Theta,\mu)$. Since $\lambda_0+\zeta=\lambda_1$,
\[
\begin{aligned}
\lambda_0\,\psi(\cdot,\Theta,\mu)+R^\psi(\Theta)&=\lambda_1\,\psi(\cdot,\Theta,\mu)+q^\psi ,
\qquad\psi\in\{b,\sigma,h,f\},\\
\lambda_0\,g(X_T,\Law(X_T))+R^g_T(\Theta)&=\lambda_1\,g(X_T,\Law(X_T))+q^g_T ,
\end{aligned}
\]
so $\Theta$ solves the forced $\lambda_1$-system \eqref{eq:forced} with forcing $q$.

Any solution of the $\lambda_1$-system is a fixed point of the same $\Phi$, which has only
one, so the solution is unique. Hence $\mathsf{WP}(\lambda_1)$, and $\zeta_0$ is
independent of $\lambda_0$ and $q$.
\end{proof}

\subsection{Proof of Theorem~\texorpdfstring{\ref{thm:cwp}}{3.5}}\label{sub:proofcwp}

\begin{proof}
By Lemma~\ref{lem:small}, $\mathsf{WP}(\lambda)$ holds for $\lambda\in[0,\lambda_*]$. With
$\zeta_0$ from Lemma~\ref{lem:cont}, set
\[
\lambda_k:=\min\{\lambda_*+k\zeta_0,\,1\},\qquad k\ge0 ,
\]
so that $\lambda_0=\lambda_*$, $\lambda_{k+1}-\lambda_k\le\zeta_0$ and $\lambda_k=1$ once
$k\ge\lceil(1-\lambda_*)/\zeta_0\rceil$. Lemma~\ref{lem:cont} carries
$\mathsf{WP}(\lambda_k)$ to $\mathsf{WP}(\lambda_{k+1})$ at every step, so
$\mathsf{WP}(1)$ holds; taking $q\equiv0$ there gives a unique solution of
\eqref{eq:mvfbsdej} in $\HH$.
\end{proof}

